\documentclass[12pt]{article}
\usepackage[T1]{fontenc}
\usepackage[utf8]{inputenc}
\usepackage{lmodern}
\usepackage{textcomp}
\usepackage{enumitem}
\usepackage{geometry}
\usepackage{graphicx}
\usepackage{amsmath, amssymb}
\geometry{margin=1in}
\begin{document}
\begin{center}
Sociology - Graduate Level Assessment 18 \
Total Marks: 40 \
Answer all questions.
\end{center}

\section*{Multiple Choice (10 marks)}
\begin{enumerate}[label=\arabic*.]
\item Frank (1967) made the claim that 'underdeveloped' societies were:
\begin{enumerate}[label=(\alph*)]
    \item insufficiently involved in the international capitalist economy
    \item reluctant to surrender their traditional ways of life
    \item economically dependent on the wealthy countries that exploited them
    \item the 'metropoles' to which 'satellite' countries were attached
\end{enumerate}
\item Pilcher (1999) identified soap operas as a 'feminine genre' of media because:
\begin{enumerate}[label=(\alph*)]
    \item most of the characters in soap operas are women
    \item they represent images of women as both domesticated and independent
    \item they alienate women and appeal to an audience of men
    \item female television producers are most likely to work in this area
\end{enumerate}
\item Sreberny-Mohammadi (1996) argues that national cultures can resist American cultural domination of the media by:
\begin{enumerate}[label=(\alph*)]
    \item domesticating its content, including more 'home-produced' programmes
    \item controlling the distribution of imported products by banning satellite dishes
    \item creating 'reverse flows' of their own programmes back to imperial societies
    \item all of the above
\end{enumerate}
\item The market model of state welfare is based on the principle of:
\begin{enumerate}[label=(\alph*)]
    \item individuals buying welfare privately, with some means-tested benefits
    \item regular benefit payments to men as earners of the 'family wage'
    \item a universalist system of welfare for all, regardless of income
    \item decommodifying social welfare through state provision
\end{enumerate}
\item In modern societies, social status is typically measured by a person's:
\begin{enumerate}[label=(\alph*)]
    \item age
    \item income
    \item verbal fluency
    \item occupation
\end{enumerate}
\end{enumerate}

\section*{True / False (10 marks)}
\textit{Answer True or False}
\begin{enumerate}[label=\arabic*.]
\setcounter{enumi}{5}
\item True or False: Ethnographic research produces qualitative data by uncovering rich, detailed accounts from an insider's perspective.
\item True or False: The Nation of Islam primarily appealed to African-Americans who felt marginalized within the broader American society.
\item True or False: Differential association theory posits that individuals learn deviant behavior through their relationships with peers.
\item True or False: The 'third age' of the life course is characterized by active non-work and independence following retirement.
\item True or False: Role-learning theory asserts that social roles are fixed and do not change over time through social interaction.
\end{enumerate}

\section*{Long Form Response (20 marks)}
\textit{Answer each question with a clear and complete explanation.}
\begin{enumerate}[label=\arabic*.]
\setcounter{enumi}{10}
\item Discuss the implications of Berger and Luckmann's concept of reality as a socially constructed phenomenon, focusing on how individuals negotiate shared definitions and the consequences of this process for social understanding.
\item Discuss the implications of Warner's study on the 'colour line' in Natchez, focusing on how it reflects the social structure and beliefs regarding racial superiority in the American Deep South.
\end{enumerate}

\vfill
\noindent\textit{End of Paper}
\end{document}